\section{Introduction}
\label{sec:introduction}

The primitive-plane route began with a four-corner coefficient in the first
mixed trace.  TPC--69 shows that the first surviving signed term is
cross-key and that its residual sign is carried by the odd erasure
boundary \citep{WangTPC69}.  TPC--70 then gives exact cofactor coordinates:
after extracting the common erasure content \(q\), the surviving sign is a
product \(\mu(a)\mu(b)\mu(u)\mu(v)\) that is frozen along each affine prime
ray \citep{WangTPC70}.  TPC--71 applies coprimality inversion and
two-dimensional Abel summation to a general finitely supported plane kernel
\citep{WangTPC71}.  Its master certificate contains
\[
 \sum_{d\le D}
 \mathfrak M_d(U/d)\mathfrak M_d(V/d)
 \cV_\square(K^{[d]})
 +
 \cT_D(K).
\]
TPC--72 carries the mass-normalized variation, the representation-dependent
diagonal tail before a completion is fixed, and complete anchor summation
as open gates; it also warns that a four-corner trace coefficient is not a
two-corner erasure entry
\citep{WangTPC72}.

The present paper re-examines only the literal support of the kernel before
trying to estimate it.  The physical coordinates are not an arbitrary pair
followed by a coprimality mask.  They are defined by
\[
 q=(S_m,S_n),\qquad S_m=qu,\qquad S_n=qv.
\]
Consequently \((u,v)=1\) identically.  If a kernel is extended by zero at
coordinates which do not occur physically, then it is already supported on
the primitive lattice.  For \(d\ge2\), the point \((dx,dy)\) cannot be
primitive.  Hence every higher diagonal slice vanishes.

This observation has two different consequences which must not be
conflated.
\begin{enumerate}[label=\textup{(\roman*)},leftmargin=3em]
 \item The support-faithful physical completion is tail-free.  Its exact
       Abel certificate has one slice and needs no growing-\(d\)
       restricted-Mertens theorem.
 \item Values away from the primitive lattice are invisible to the
       physical signed sum.  They may therefore be introduced as an
       analytic completion.  Such a completion can fill hard coprimality
       holes and reduce variation, but its higher diagonal slices and tail
       return and must be charged.
\end{enumerate}
Thus variation and diagonal tail are not independent invariants of the
primitive physical datum.  They form a coupled analytic completion gauge.
Tail freedom is canonical for the support-faithful zero extension; low
variation is a separate theorem.

\subsection{Main results}

The paper proves seven groups of statements.
\begin{enumerate}[label=\textup{(\arabic*)},leftmargin=3em]
 \item The exact \(q\)-coordinates make every literal anchored kernel
       primitive-supported after the \(q\)-sum and the removal of the
       same-residual direction.
 \item The canonical zero completion annihilates all \(d\ge2\) diagonal
       slices and every exact diagonal tail with cutoff \(D\ge1\).
 \item The anchored coefficient and the complete scalar mixed trace have
       exact \(d=1\) Abel identities and weighted-plaquette majorants.
 \item A Banach-valued, anchor-labelled version retains the complete
       growing anchor family and replaces uniform anchorwise little-oh
       notation by one global cost inequality.
 \item Every finitely supported kernel has a unique terminal-rectangle
       expansion.  Mixed variation is precisely the \(\ell^1\)-mass of
       its rectangle charges, and diagonal downsampling is contractive.
 \item Every analytic completion of the same primitive datum gives the
       same physical sum.  Taking the infimum of the full divisor-variation
       cost over declared completions is an exact sufficient certificate.
 \item Counterexamples show that zero completion need not be
       variation-optimal, completion costs can be inflated arbitrarily,
       primitive support need not cancel, and scalar trace cancellation
       does not control the erasure operator.
\end{enumerate}

\subsection{Logical level}

The finite-support identities and completion calculus are L0.  The
identification of the literal fixed-\(h_0\) kernel with its canonical
primitive zero completion is L1: it retains all complete native keys,
actual support, literal coefficients, and raw atomic normalization.  The
paper does not estimate the resulting physical weighted plaquette cost
relative to the anchored or trace mass.  That missing inequality is the
next L2 gate.  In particular, eliminating an auxiliary diagonal tail is
not a parity breakthrough.

\subsection{Organization}

\Cref{sec:physical-support} fixes the physical datum and its support.
\Cref{sec:tail-collapse} proves tail collapse and the single-slice Abel
form.  \Cref{sec:rectangle-charges} gives the canonical rectangle
decomposition.  \Cref{sec:completion-gauge} develops analytic completions.
\Cref{sec:q-aggregation} treats the order of \(q\), anchor, and trace
aggregation.  Countermodels appear in \cref{sec:countermodels}.
\Cref{sec:physical-ledger} states the exact remaining physical gate, and
\cref{sec:route-plan,sec:claim-boundary} give the adaptive next route and
the final claim ledger.
